\documentclass[11pt]{article}
\usepackage{fontspec}
\usepackage[french]{babel}
\usepackage{amsmath}
\usepackage{amssymb}
\usepackage{amsthm}
\usepackage{mathrsfs}
\usepackage{tikz-cd}
\usepackage{enumitem}

\newcommand{\printedpage}[1]{\marginpar{\footnotesize p.~#1}}
\newcommand{\msp}[1]{\marginpar{\footnotesize [#1]}}
\newcommand{\cercle}[1]{\tikz[baseline=(X.base)]\node[draw,circle,inner sep=1pt](X){\upshape #1};}
\newcommand{\defeq}{\overset{\text{\upshape déf}}{=}}
\newcommand{\cardred}{\operatorname{card}\,\text{\upshape réd}\,}
\newcommand{\Hom}{\underline{\operatorname{Hom}}}
\newcommand{\uGamma}{\underline{\Gamma}}
\newcommand{\Ind}{\operatorname{Ind}}
\newcommand{\into}{\lhook\joinrel\longrightarrow}

\begin{document}

\printedpage{99}
% head: A. Grothendieck : Dérivateurs, Ch. XVIII, Édition des Universités de Montpellier II et Paris VII
% foot: Version 17 décembre 2025

un homomorphisme sur les $\varinjlim$ (\footnote{Attention, c'est ici qu'on utilise tacitement le fait que le foncteur $\lambda_* : \mathcal{E}_{j'} \longrightarrow \mathcal{E}_j$ commute aux $\varinjlim$ de type $I_\alpha$, ce qui implique que les inclusions $\mathcal{E}_j \into \mathcal{E}$ commutent à ce type de $\varinjlim$ .}). La condition (i) assure que $\lambda \longmapsto G(\lambda)$ transforme identités en identités, et la condition (iii) que l'on a $G(\mu\lambda) = G(\mu)G(\lambda)$. Donc $G$ apparaît comme une section de $\mathcal{E}$ sur $J$. On a d'autre part des flèches
\[
X(j) \defeq X_0(j) \longrightarrow G(j) = \varinjlim_{\beta\in I_\alpha} X_\beta(j) \ ,
\]
ce qui prouve F (1). D'autre part, la condition 3°) du lemme prouve la validité de F$'$ (1). Cela achève la démonstration de 4.10.3, et par là de 4.10.1 et de 4.9.3.

\msp{page 263}

\textbf{Corollaire 4.10.8} (du théorème 4.10.1). \textit{Soit $\mathcal{N} \longrightarrow J$ un foncteur cofibrant, tel que les fibres soient des catégories accessibles (4.8), les foncteurs cochangement de base $\lambda_*$ : $\mathcal{N}_{j'} \longrightarrow \mathcal{N}_j$ accessibles, enfin $J$ équivalente à une petite catégorie. Alors la catégorie $\mathcal{M} = \uGamma(\mathcal{N}/J)$ est accessible. Plus précisément, soit $\pi_0$ un cardinal infini tel que les catégories $\mathcal{N}_j$ soient $\pi_0$-accessibles, et que les foncteurs $\lambda_*$ soient $\pi_0$-accessibles, et soit $\pi$ un cardinal infini tel que l'on ait :}
\begin{enumerate}[leftmargin=*]
\item[\cercle{a}] $\pi^{\pi_0} = \pi$.
\item[\cercle{b}] $\forall\ \lambda : j' \longrightarrow j$ \textit{dans} $J$, $\lambda_* : \mathcal{E}_j \longrightarrow \mathcal{E}_{j'}$ \textit{applique} $(\mathcal{E}_j)_{\pi_0}$ \textit{dans} $(\mathcal{E}_{j'})_{\pi_0}$.
\item[\cercle{c}] $\cardred J \leqslant \pi$.
\end{enumerate}
\textit{Alors $\mathcal{M}$ est $\pi$-accessible, et plus précisément, on a}
\[
\mathcal{M} \ \simeq \ \Ind_\pi(\mathcal{M}_\pi) \ ,
\]
\textit{où $\mathcal{M}_\pi$ est la sous-catégorie pleine de $\mathcal{M}$ formée des sections $F$ de $\mathcal{M}$ sur $J$ telles que l'on ait $F(j) \in (\mathcal{E}_j)_\pi \ \forall\ j \in \operatorname{Ob} J$,} \msp{page 264} \textit{$(\mathcal{E}_j)_\pi$ désignant la sous-catégorie pleine de $\mathcal{E}_j$ formée des objets $\pi$-accessibles de $\mathcal{E}_j$.}

En effet, l'hypothèse au début implique l'existence de $\pi_0$ et de $\pi$ satisfaisant les conditions dites. Les conditions a) et b) sur $\pi$, plus le fait que les foncteurs $\lambda_*$ soient $\pi_0$-accessibles, impliquent que la sous-catégorie strictement pleine $\mathcal{N}_\pi$ de $\mathcal{N}$ dont les fibres sont les $(\mathcal{N}_j)_\pi$, est une sous-catégorie cofibrée, et que de plus $\mathcal{N}$ est $J$-équivalent à $\Ind_\pi(\mathcal{N}_\pi/J)$. D'autre part, les fibres de $\mathcal{N}_\pi/J$ sont stables par $\varinjlim$ grandes devant $\pi_0$ et de cardinal $\leqslant \pi$, donc la condition “ordinale” de 4.10.1 est vérifiée (en y faisant $\mathcal{D} = \mathcal{N}_\pi$). Comme de plus $\operatorname{card} J \leqslant \pi$, on en conclut par 4.10.1 que $\mathcal{M} \simeq \Ind_\pi(\mathcal{M}_\pi)$.

\msp{page 265}

\textbf{Corollaire 4.10.9.} \textit{Soit $\mathcal{N} \longrightarrow J$ une “cofibration en paratopos” (\footnote{terme ambigu, cf. la double explication.}), i.e. un foncteur cofibrant tel que les fibres soient des catégories paratopos, et les foncteurs cochangement de base soient des foncteurs images inverses (resp. images directes) pour des morphismes de paratopos, i.e. qu'ils commutent aux (petites) $\varinjlim$ (resp. aux petites $\varprojlim$ ). Supposons de plus $J$ équivalente à une petite catégorie. Alors $\uGamma(\mathcal{N}/J)$ est une catégorie paratopos.}

En effet, c'est une catégorie accessible en vertu de 4.10.8, d'autre part l'hypothèse sur les $\lambda_*$ implique aussitôt que $\uGamma(\mathcal{N}/J)$ est stable par petites $\varprojlim$ (resp. par petites $\varprojlim$
\printedpage{100}
% head: A. Grothendieck : Dérivateurs, Ch. XVIII, Édition des Universités de Montpellier II et Paris VII
% foot: Version 17 décembre 2025
\texttt{[plutôt} $\varinjlim$ \texttt{?]}). Donc $\uGamma(\mathcal{N}/J)$ est une catégorie paratopos en vertu de 4.8.9 b $\Longrightarrow$ a (resp. c $\Longrightarrow$ a).

\textbf{Corollaire 4.10.10.} \textit{Soit $\mathcal{N} \longrightarrow J$ une catégorie} fibrée \textit{(non cofibrée!), telle que les fibres soient des catégories paratopos et les foncteurs} \msp{page 266} \textit{changement de base soient des foncteurs images directes pour des morphismes de paratopos (i.e. des foncteurs accessibles, commutant aux petites $\varprojlim$ ). Supposons de plus $J$ équivalente à une petite catégorie. Alors $\uGamma(\mathcal{N}/J)$ est une catégorie paratopos.}

En effet, l'hypothèse implique que pour toute flèche $\lambda : j' \longrightarrow j$ dans $J$, le foncteur $\lambda_*$ : $\mathcal{N}_{j'} \longrightarrow \mathcal{N}_j$ admet un adjoint à gauche $\lambda^* : \mathcal{N}_j \longrightarrow \mathcal{N}_{j'}$, d'où s'ensuit que $\mathcal{N} \longrightarrow J$ est aussi cofibrant, les foncteurs de cochangement de base commutant aux petites $\varinjlim$ . Donc on peut appliquer 4.8.9, qui donne le résultat voulu.

\textbf{Question 4.10.11.} \textit{Soit $\mathcal{N} \longrightarrow J$ un foncteur} fibrant, \textit{tel que les $\mathcal{N}_j$ soient des catégories accessibles, et les foncteurs $\lambda^*$ de changement de base accessibles, $J$ étant équivalente à une petite catégorie. Je} \msp{page 267} \textit{conjecture que sous ces conditions, $\uGamma(\mathcal{N}/J) \defeq \Hom(J,\mathcal{N})$ est accessible, et qu'il en est de même de $\Hom\,\underline{\operatorname{cart}}_J(J,\mathcal{N})$.}

C'est ce qui est prouvé dans SGA 4 I 9.24 sous des hypothèses légèrement plus fortes (à l'aide de l'ancienne notion trop forte à mon gré, de filtrations cardinales). Je serais étonné que la démonstration (un peu longuette) de loc. cit. ne s'applique pas au cas général. Mais on aimerait avoir un énoncé qui coiffe celui-ci et 4.10.8, dans le cas d'un foncteur $\mathcal{N} \longrightarrow J$ qui ne serait supposé ni fibrant ni cofibrant, mais \ldots(???).

\section*{4.11\quad Application à l'existence des limites dans une catégorie de la forme $\Ind_{\pi_0}(C)$.}

\textbf{Proposition 4.11.1.} \textit{Soit} \msp{page 268} \textit{$C$ une catégorie équivalente à une petite catégorie, $\pi$ un cardinal infini, posons $\mathcal{M} = \Ind_\pi(C)$. Soit $\Phi$ un ensemble de catégories de cardinal réduit $\leqslant \pi$. On se propose d'examiner si $\mathcal{M}$ est “stable par $\varinjlim$ de type $\Phi$” (i.e. stable par $\varinjlim$ de type $I$, quel que soit $I \in \Phi$). Il suffit pour cela que $C$ est stable par $\varinjlim$ de type $\Phi$, et que de plus il existe un ordinal limite $\alpha$ avec $\operatorname{card}\alpha \leqslant \pi$ et tel que $C$ soit stable par $\varinjlim$ de type $I_\alpha$ (p.ex. tel que $I_\alpha \in \Phi$). Inversement, si $\mathcal{M}$ est stable par $\varinjlim$ de type $\Phi$, et si on suppose de plus $C$ karoubienne (i.e. stable par facteurs directs), alors $C$ est stable par $\varinjlim$ de type $\Phi$. Dans l'un et l'autre cas, le foncteur d'inclusion $C \into \mathcal{M}$ commute aux $\varinjlim$ envisagées.}

(\footnote{Expliciter ici les corollaires des types 4.13.20, 4.13.21 (p. 341).}).

\textsc{Démonstration}. a) Condition suffisante. On applique 4.9.3 pour conclure que
\printedpage{101}
% head: A. Grothendieck : Dérivateurs, Ch. XVIII, Édition des Universités de Montpellier II et Paris VII
% foot: Version 17 décembre 2025
\msp{page 269}
pour toute catégorie $J$ de cardinal réduit $\leqslant \pi$ (en particulier pour $J \in \Phi$), tout foncteur $\varphi : J \longrightarrow \mathcal{M} = \Ind_\pi(C)$ admet une représentation indicielle $\Phi : J \times I \longrightarrow C$, avec $I$ grande devant $\pi$. Ceci, plus la stabilité de $C$ par $\varinjlim$ de type $J$, suffit à établir l'existence de la $\varinjlim$ de $\varphi$.

La réciproque provient du fait que si $C$ est karoubienne, alors son image essentielle dans $\mathcal{M}$ n'est autre que $\mathcal{M}_\pi$, laquelle est stable dans $\mathcal{M}$ par toute $\varinjlim$ de cardinal $\leqslant \pi$ représentable dans $\mathcal{M}$. Si $\mathcal{M}$ est stable par $\varinjlim$ de type $\Phi$, il en est de même de $\mathcal{M}_\pi$, donc de $C$, q.e.d.

\textbf{Corollaire 4.11.2.} \textit{Pour que $\mathcal{M}$ soit stable par $\varinjlim$ filtrantes (resp. par petites $\varinjlim$ ), il suffit que $C$ le soit par $\varinjlim$ filtrantes (resp. par $\varinjlim$ )} de cardinal $\leqslant \pi$, \textit{et c'est aussi nécessaire si $C$ est karoubienne.}

Supposons $C$ stable par $\varinjlim$ filtrantes \msp{page 270} (resp. par $\varinjlim$ de cardinal $\leqslant \pi$, alors par 4.11.1 $\mathcal{M}$ est stable par ce même type de limites inductives. D'autre part, il est stable par $\varinjlim$ grandes devant $\pi$. Il suffit donc de prouver le

\textbf{Lemme 4.11.3.} \textit{Soit $\mathcal{M}$ une catégorie, $\pi$ un cardinal. Supposons que $\mathcal{M}$ soit stable par $\varinjlim$ filtrantes de cardinal $\leqslant \pi$ (resp. par $\varinjlim$ de cardinal $\leqslant \pi$, resp. par sommes des cardinal $\leqslant \pi$),} et \textit{par $\varinjlim$ grandes devant $\pi$. Alors $\mathcal{M}$ est stable par petites $\varinjlim$ filtrantes (resp. par petites $\varinjlim$ , resp. par petites sommes).} % [sic] des cardinal -> de cardinal, agreement (cf. "de cardinal" elsewhere in the statement)

Ce lemme se prouve comme 4.4.14, page 97.

Reste la réciproque dans 4.11.2, qui résulte de l'assertion similaire dans 4.11.1.

\textbf{Corollaire 4.11.4.} \textit{Quand $\mathcal{M} = \Ind_\pi C$ est-il stable par sommes ? S'il existe un ordinal $\alpha$ tel que $\operatorname{card}\alpha \leqslant \pi$, et que $C$ soit stable par $\varinjlim$ de type $I_\alpha$, alors il suffit que $C$ soit stable par sommes de cardinal $\leqslant \pi$. Cette condition est aussi nécessaire si $C$ est karoubienne.}

(\footnote{Pour autres variantes de 4.11.1, cf. 4.13.19 (p. 340) et corollaires.}).

Se prouve comme 4.11.2, en utilisant 4.11.3.

\msp{page 271}

\textbf{Proposition 4.11.5} (\footnote{?, démonstration pas achevée.}). \textit{Soit $\Phi$ un ensemble de catégories $I$ équivalentes à des catégories finies, et “rigides”, i.e. pour tout $i \in \operatorname{Ob} I$, $\Hom_I(i,i) = \{\mathrm{id}_i\}$. Soient $C$ une catégorie équivalente à une petite catégorie, $\pi$ un cardinal infini, $\mathcal{M} = \Ind_\pi C$. Pour que $\mathcal{M}$ soit stable par $\varinjlim$ (resp. par $\varprojlim$ ) de type $\Phi$, il} suffit \textit{que $C$ le soit. Dans le cas non respé, si $C$ est karoubienne, cette condition est aussi nécessaire.}

Cette dernière assertion est contenue dans l'assertion similaire dans 4.11.1. Pour la première assertion, cela résulte de la façon habituelle du

\textbf{Lemme 4.11.6} (\footnote{?, démonstration non explicitée est peut-être faux!}). \textit{Soit $J$ une catégorie équivalente à une catégorie finie, $J$ rigide, $C$ une catégorie équivalente à une petite catégorie, $\pi$ un cardinal infini. Alors tout foncteur}
\printedpage{102}
% head: A. Grothendieck : Dérivateurs, Ch. XVIII, Édition des Universités de Montpellier II et Paris VII
% foot: Version 17 décembre 2025
\textit{$J \longrightarrow \mathcal{M} \defeq \Ind_\pi C$ admet une représentation indicielle $J \times I \longrightarrow C$, avec $I$ grande devant $\pi$, en d'autres termes, le} \msp{page 272} \textit{foncteur canonique pleinement fidèle}
\[
\Ind_\pi(\Hom(J,C)) \ \longrightarrow \ \Hom(J,\mathcal{M})
\]
\textit{déduit de l'inclusion $C \into \mathcal{M}$, d'où $\Hom(J,C) \into \Hom(J,\mathcal{M})$ (cf. 4.9), est une équivalence de catégories.} (\textbf{N.B.} J'ignore s'il en est ainsi sans hypothèse de rigidité sur $J$.)

\textsc{Démonstration}. Cela se prouve comme le résultat similaire pour $\Ind(C)$ au lieu de $\Ind_\pi C$, cf. SGA 4 I 8.8.5.

\msp{page 273}

% a vertical rule is printed in the left margin beside Corollaire 4.11.7
\textbf{Corollaire 4.11.7} (\footnote{sous réserve de validité de 4.11.5.}). \textit{Pour que $\mathcal{M} = \Ind_\pi C$ soit stable par $\varinjlim$ (resp. par $\varprojlim$ )} finies, \textit{il} suffit \textit{que $C$ le soit. Et c'est aussi nécessaire dans le cas non respé, lorsqu'on suppose $C$ karoubienne.}

\section*{4.12\quad Foncteurs accessibles entre catégories accessibles}

\textbf{Proposition 4.12.1} (\footnote{Remonterait à 4.8.}). \textit{Soient $\mathcal{M}$, $\mathcal{N}$ deux catégories accessibles (4.8), et $f : \mathcal{M} \longrightarrow \mathcal{N}$ un foncteur accessible. Alors il existe des cardinaux infinis $\pi_0$, $\pi_1$ tels que $\pi_1 \geqslant \pi_0$ et que pour tout cardinal $\pi \geqslant \pi_1$, on ait}
\[
f(\mathcal{M}_\pi) \ \subset \ \mathcal{N}_{\pi^{\pi_0}} \ , \tag*{(4.12.1)} % [sic] (4.12.1) -> (4.12.1.1), the formula is cited below as (4.12.1.1) and the next one is numbered (4.12.1.2)
\]
\textit{et en particulier}
\[
f(\mathcal{M}_\pi) \ \subset \ \mathcal{N}_\pi \qquad \textit{si } \pi = \pi^{\pi_0}. \tag*{(4.12.1.2)}
\]
\textit{De façon plus précise, on peut prendre pour $\pi_0$ un cardinal utile pour $\mathcal{M}$} \msp{page 274} \textit{et tel que $f$ soit $\pi_0$-accessible (p.ex. le plus petit cardinal de ce type). Soit alors $\pi_1$ un cardinal tel qu'on ait $f(\mathcal{M}_{\pi_0}) \subset \mathcal{N}_{\pi_1}$, $\pi_1 \geqslant \cardred \mathcal{M}_{\pi_0}$, et enfin tel que $\mathcal{N}_{\pi_1}$ soit} \texttt{[une]} \textit{sous-catégorie strictement génératrice de $\mathcal{N}$. Le couple $(\pi_0,\pi_1)$ satisfait aux conditions énoncées.}

\textsc{Démonstration}. On aura en effet
\[
\mathcal{N}_1 \ \xleftarrow{\ \sim\ } \Ind_{\pi_0}(\mathcal{M}_{\pi_0})
\]
et $f$ commutant aux $\varinjlim$ grandes devant $\pi_0$, sera connu (à isomorphisme unique près) par sa restriction à $\mathcal{M}_{\pi_0}$, laquelle se factorise par $\mathcal{N}_{\pi_1}$,
\[
f|\mathcal{M}_{\pi_0} : \mathcal{M}_{\pi_0} \ \longrightarrow \ \mathcal{N}_{\pi_1} \ .
\]
Pour prouver (4.12.1.1), posant $\pi' = \pi^{\pi_0} \geqslant \pi$, et notant que $\pi'^{\pi_0} = \pi_0$ \texttt{[plutôt} $\pi'^{\pi_0} = \pi'$\texttt{]},
\[
f(\mathcal{M}_\pi) \ \subset \ f(\mathcal{M}_{\pi'}) \ \overset{?}{\subset} \ \mathcal{N}_{\pi'} \ ,
\]
\printedpage{103}
% head: A. Grothendieck : Dérivateurs, Ch. XVIII, Édition des Universités de Montpellier II et Paris VII
% foot: Version 17 décembre 2025
on voit qu'il suffit de prouver (4.12.1.2) (pour $\pi'$). Donc on peut supposer $\pi = \pi^{\pi_0}$ et prouver (4.12.1.2). Mais comme $\pi_1 \geqslant \cardred \mathcal{M}_{\pi_0}$, on sait que si $\pi$ est tel que $\pi \geqslant \pi_1$, $\pi^{\pi_0} = \pi$, alors $\mathcal{M}_\pi$ est formé des $\varinjlim\limits_{I} x_i$, les $x_i \in \mathcal{N}_{\pi_0}$, et $I$ grande devant $\pi_0$, et $\operatorname{card} I \leqslant \pi$.

\msp{page 275}

On a alors ($f$ étant $\pi_0$-accessible)
\[
f(x) \ = \ \varinjlim_{I} f(x_i) \ ,
\]
les $f(x_i)$ sont dans $\mathcal{N}_{\pi_1}$, donc dans $\mathcal{N}_\pi$, or $\mathcal{N}_\pi$ stable par $\varinjlim$ de cardinal $\leqslant \pi$ (puisque $\mathcal{N}$ est engendrée par $\mathcal{N}_{\pi_1}$, formée d'objets $\pi_1$-accessibles et a fortiori $\pi$-accessibles, cf. SGA 4 I 9.9). D'où $f(x) \in \mathcal{N}_\pi$, q.e.d.

\textbf{Corollaire 4.12.2.} \textit{Supposons que $\pi_0$ soit aussi utile pour $\mathcal{N}$, et que $\pi_1 \geqslant \cardred \mathcal{N}_{\pi_0}$. Soit $\pi$ un cardinal $\geqslant \pi_1$, et tel que $\pi^{\pi_0} = \pi$. Alors $f$ s'identifie, à équivalence de catégories et à isomorphisme de foncteurs près, au foncteur}
\[
\Ind(f_\pi) \ : \ \Ind_\pi(\mathcal{M}_\pi) \ \longrightarrow \ \Ind_\pi(\mathcal{N}_\pi) \ ,
\]
\textit{où $f_\pi : \mathcal{M}_\pi \longrightarrow \mathcal{N}_\pi$ est le foncteur induit par $f$.}

\textbf{Proposition 4.12.3} (\footnote{remonte à 4.7.}). \textit{Soient $C$ une catégorie équivalente à une petite catégorie, $\mathcal{N}$ une $\mathfrak{U}$-catégorie, $\pi_0$ un cardinal infini.} \msp{page 276} \textit{On suppose que $\mathcal{N}$ satisfait $L_{\pi_0}$ (stabilité par $\varinjlim$ grandes devant $\pi_0$). Soit $\mathcal{M} = \Ind_{\pi_0} C$, $\Hom_{!\,\pi_0}(\mathcal{M},\mathcal{N})$ la sous-catégorie strictement pleine de $\Hom(\mathcal{M},\mathcal{N})$ formée des $f$ : $\mathcal{M} \longrightarrow \mathcal{N}$ $\pi_0$-accessibles. Alors le foncteur restriction}
\[
\Hom_{!\,\pi_0}(\mathcal{M},\mathcal{N}) \ \longrightarrow \ \Hom(C,\mathcal{N}) \tag*{(4.12.3.1)}
\]
\textit{est une équivalence de catégories. On trouve un foncteur quasi-inverse comme composé}
\[
\Hom(C,\mathcal{N}) \ \longrightarrow \ \Hom(\underbrace{\Ind_{\pi_0}(C)}_{\mathcal{M}},\Ind_{\pi_0}(\mathcal{N})) \ \longrightarrow \ \Hom(\mathcal{M},\mathcal{N}) \ , \tag*{(4.12.3.2)}
\]
\textit{où la deuxième flèche est le foncteur composition $F \longmapsto L \circ F$ avec le foncteur limite inductive $L : \Ind_{\pi_0}\mathcal{N} \longrightarrow \mathcal{N}$. Le foncteur composé (4.11.3.2) a une image contenue dans $\Hom_{!\,\pi_0}(\mathcal{M},\mathcal{N})$ et induit une équivalence entre cette dernière catégorie et $\Hom(C,\mathcal{N})$.} % [sic] (4.11.3.2) -> (4.12.3.2), the composed functor is the display just labelled (4.12.3.2)

\msp{page 277}

\textit{En d'autres termes, on peut dire que le foncteur d'inclusion canonique}
\[
C \ \longrightarrow \ \Ind_{\pi_0} C \ = \ \mathcal{M}
\]
\textit{est} 2-universel \textit{pour les foncteurs de $C$ dans des catégories $\mathcal{N}$ satisfaisant $L_{\pi_0}$, celles-ci étant considérées comme $0$-objets d'une $2$-catégorie où les $1$-objets sont les foncteurs $\mathcal{N} \longrightarrow \mathcal{N}'$ $\pi_0$-accessibles, i.e. commutent aux $\varinjlim$ grandes devant $\pi_0$.}

Cet énoncé est l'analogue, pour $C$ petite quelconque, de 4.4.16 (page 101), valable pour $C$ stable par $\varinjlim$ de cardinal $\leqslant \pi_0$. Il y a bien sûr d'autres variantes, en relation avec la commutation à d'autres types de $\varinjlim$ ( $\varinjlim$ filtrantes, sommes, limites inductives finies etc.) envisagées dans la section 4.11 précédente.

La démonstration de 4.12.3 est soritale.

\end{document}
